\documentclass[journal,twocolumn,10pt]{IEEEtran}

\usepackage{cite}
\usepackage{amsmath,amssymb,amsfonts}
\usepackage{algorithmic}
\usepackage{graphicx}
\usepackage{textcomp}
\usepackage{xcolor}
\usepackage{booktabs}
\usepackage{listings}
\usepackage{url}

\lstset{
  basicstyle=\ttfamily\footnotesize,
  breaklines=true,
  frame=single,
  numbers=none,
  keywordstyle=\color{blue},
  commentstyle=\color{gray}
}

\begin{document}

\title{SiliconTree: An Interactive, Cycle-Accurate Microarchitectural Simulation and Speculative Vulnerability Analysis Suite}

\author{Sridhar~Shanmugam,~\IEEEmembership{Member,~IEEE}
\thanks{Sridhar Shanmugam is with the Department of Computer Science and Engineering, SiliconTree Architecture Research Group (e-mail: sridhar@silicontree.org).}}

\markboth{IEEE TRANSACTIONS ON CIRCUITS AND SYSTEMS---II: EXPRESS BRIEFS, VOL.~72, NO.~3, MARCH~2026}%
{Shanmugam: SiliconTree: Interactive Cycle-Accurate Microarchitectural Simulation Suite}

\maketitle

\begin{abstract}
Modern microprocessors employ aggressive performance optimizations, including deep pipelining, dynamic branch prediction, out-of-order execution, and multi-level caching. However, these complex microarchitectural optimizations introduce severe hardware security vulnerabilities, most notably transient and speculative execution side-channel attacks such as Spectre and Meltdown. Existing architectural simulators suffer from an intractable trade-off: heavyweight industrial frameworks (e.g., gem5) exhibit steep learning curves and lack accessible real-time microarchitectural state observability, whereas lightweight educational simulators (e.g., MARS, RARS) omit pipeline hazards, cache conflict dynamics, and speculative execution side effects altogether. In this paper, we present SiliconTree, an interactive, cycle-accurate computer architecture and hardware security simulation suite built with strict Model-View-Controller (MVC) decoupling. SiliconTree integrates a cycle-accurate 5-stage RISC processor featuring reverse-order pipeline evaluation ($\text{WB} \to \text{MEM} \to \text{EX} \to \text{ID} \to \text{IF}$), Radix-2 Booth and Radix-4 Bit-Pair recoding multipliers, restoring and non-restoring integer dividers, $N$-way set-associative LRU caches with write-back/dirty tracking, and bus arbitration with cycle-stealing DMA. Crucially, SiliconTree implements microarchitectural transient execution tracking, faithfully capturing out-of-bounds speculative cache line warming, Flush+Reload timing recovery, and $\texttt{LFENCE}$ serialization barriers. SiliconTree records immutable \texttt{CycleSnapshot} states into an in-memory deque, enabling instantaneous single-cycle time-travel ($-1$ Cyc) and post-mortem forensic timeline scrubbing. Integrated with an inline 60 FPS SVG datapath visualization web workbench, SiliconTree demonstrates cycle-accurate hazard resolution, microarchitectural side-channel leakage detection, and architectural design-space exploration.
\end{abstract}

\begin{IEEEkeywords}
Computer architecture, microarchitectural attacks, hardware security, speculative execution, Spectre v1, pipeline hazards, cycle-accurate simulation, cache side-channel, Flush+Reload.
\end{IEEEkeywords}

\section{Introduction}
\IEEEPARstart{T}{he} relentless demand for computational throughput has driven microprocessor designers to implement sophisticated microarchitectural techniques, including superscalar pipelines, dynamic two-bit branch prediction, out-of-order execution, and multi-level hierarchical caching \cite{hennessy2017, hill2008}. While these mechanisms achieve substantial instructions-per-cycle (IPC) gains, they introduce microarchitectural side channels \cite{ge2021, chandrasekaran2022}. As demonstrated by transient execution attacks such as Spectre \cite{kocher2019} and Meltdown \cite{lipp2018}, microarchitectural state changes (e.g., lines warmed in the L1 data cache during speculatively executed, architecturally squashed code paths) fail to roll back when a branch misprediction is resolved. Consequently, sensitive data can be exfiltrated across privilege boundaries via cache timing discrepancies such as Flush+Reload \cite{yarom2014}.

Despite the profound security implications of hardware vulnerabilities, understanding and analyzing these microarchitectural mechanisms remains challenging \cite{zombieload2019, rowhammer2016}. Researchers and educators rely heavily on simulation tools, which predominantly fall into two extreme paradigms:
\begin{enumerate}
    \item \textbf{Full-System Architectural Emulators:} Tools such as gem5 \cite{binkert2011} provide high architectural accuracy and full-system execution. However, they incur massive computational overhead, produce multi-gigabyte text traces, require steep C++/Python domain knowledge, and lack real-time visual inspection of transient pipeline latches and cache tag dynamics.
    \item \textbf{Instruction-Level Academic Emulators:} Tools such as MARS \cite{vollmar2006} and RARS \cite{rars} offer accessible pedagogical environments but model only programmer-visible instruction set architecture (ISA) state. They omit pipeline hazard bubbles, forwarding bypass multiplexers, set-associative cache conflict misses, bus arbitration, and speculative transient execution.
\end{enumerate}

To bridge this critical divide between pedagogical accessibility and cycle-accurate security fidelity, this paper introduces \textbf{SiliconTree}, an interactive, cycle-accurate hardware simulation suite. 

The primary contributions of this paper are as follows:
\begin{itemize}
    \item \textbf{Strict MVC Hardware Simulation Core:} We design a pure-Python, zero-GUI-dependency cycle-accurate simulation engine (\texttt{core\_engine/}) implementing a 32-bit RISC ISA, dynamic branch prediction, and multi-level LRU caching.
    \item \textbf{Deterministic Reverse-Order Pipeline Evaluation:} We formulate and implement an execution paradigm that steps pipeline stages in reverse ($\text{WB} \to \text{MEM} \to \text{EX} \to \text{ID} \to \text{IF}$) on every clock tick, deterministically preventing intra-cycle latch overwrites without complex double-buffering.
    \item \textbf{Microarchitectural Vulnerability \& Barrier Modeling:} SiliconTree provides cycle-accurate modeling of Spectre Variant 1 (Bounds Check Bypass), tracking speculative out-of-bounds cache warming, Flush+Reload cycle timing extraction, and $\texttt{LFENCE}$ execution serialization.
    \item \textbf{Immutable Time-Travel \& Forensic Timeline Scrubbing:} We implement an immutable snapshot architecture (\texttt{CycleSnapshot}) stored in a state deque, providing instantaneous single-cycle stepping, rewinding ($-1$ Cyc), and crash-dump root-cause isolation.
    \item \textbf{Interactive Multi-Mode Exploration:} We validate SiliconTree across five distinct architectural and hardware security modes (SpectreScope, Silicon Heist, Silicon Deathmatch, Silicon Forensic, and Architect Tycoon) backed by an inline-SVG Flask Web Workbench.
\end{itemize}

\section{Background \& Architectural Foundations}

\subsection{Pipelining, Data Hazards, and Forwarding}
A canonical 5-stage RISC processor divides instruction execution into Instruction Fetch (IF), Instruction Decode/Register Read (ID), Execute (EX), Memory Access (MEM), and Write-Back (WB) \cite{hennessy2017}. Pipelining introduces structural, data, and control hazards:
\begin{itemize}
    \item \textbf{Read-After-Write (RAW) Hazards:} Occur when instruction $I_j$ consumes register $R_s$ before preceding instruction $I_i$ writes to $R_d$. SiliconTree implements hazard detection and forwarding bypasses directly from the $\text{EX/MEM}$ and $\text{MEM/WB}$ pipeline latches into the $\text{ID/EX}$ stage inputs, eliminating stalls for ALU-to-ALU dependencies.
    \item \textbf{Load-Use Hazards:} Occur when an instruction immediately following a load ($\texttt{LW}$) consumes the loaded data. Because memory data is only available at the conclusion of $\text{MEM}$, forwarding cannot travel backward in time; SiliconTree automatically injects a single-cycle hardware bubble ($\texttt{NOP}$) into $\text{ID}$ while holding $\text{IF}$ stalled.
\end{itemize}

\subsection{Dynamic Branch Prediction and Speculative Execution}
To prevent multi-cycle stalls on conditional branch instructions ($\texttt{BEQ}, \texttt{BNE}, \texttt{BLT}, \texttt{BGE}$), modern processors implement dynamic branch predictors. SiliconTree incorporates a 2-Bit Saturating Counter state machine coupled with a Branch Target Buffer (BTB). The predictor transitions across four states: Strongly Not Taken ($00_2$), Weakly Not Taken ($01_2$), Weakly Taken ($10_2$), and Strongly Taken ($11_2$).

When a conditional branch enters $\text{IF}$, the predictor speculates on the branch outcome. If predicted taken, the program counter ($PC$) is updated to the target address stored in the BTB; if predicted not taken, $PC$ advances sequentially. The branch condition and target are definitively resolved in $\text{EX}$. Upon misprediction detection, a penalty of 2 stall cycles is incurred: the processor squashes speculative instructions currently residing in $\text{IF}$ and $\text{ID}$, resetting stage latches to $\texttt{NOP}$ bubbles, and redirects $PC$ to the correct path.

\subsection{Speculative Side-Channel Exploits (Spectre v1)}
The Spectre Variant 1 vulnerability exploits conditional branch misprediction to bypass bounds checks \cite{kocher2019}. Consider the classic vulnerable gadget:
\begin{lstlisting}[language=C]
if (x < array1_size) {
    y = array2[array1[x] * 64];
}
\end{lstlisting}

Under normal conditions, accessing an index $x \ge \text{array1\_size}$ throws an exception or halts. However, if the branch predictor is trained to predict the branch as \textit{taken}, the CPU speculatively executes the body before the branch condition evaluates in $\text{EX}$:
\begin{enumerate}
    \item The out-of-bounds byte $k = \text{array1}[x]$ is speculatively loaded into a transient register.
    \item A secondary dependent load accesses $\text{array2}[k \times 64]$. This address maps into a specific L1 cache line, warming that cache set from DRAM.
    \item In cycle $t_{\text{resolve}}$, the branch condition evaluates in $\text{EX}$, detects the bounds violation, and rolls back architectural register state.
\end{enumerate}

Crucially, \textbf{the microarchitectural cache tag state is not rolled back}. The adversary subsequently probes every set of $\text{array2}$ using the Flush+Reload technique \cite{yarom2014}, measuring memory access latency:
\begin{equation}
\text{Latency}(\text{Set}_i) = \begin{cases} 
t_{\text{hit}} \approx 1 \text{ cycle}, & \text{if } i = k \\
t_{\text{miss}} \approx 8\text{--}10 \text{ cycles}, & \text{if } i \neq k 
\end{cases}
\end{equation}
The set exhibiting the lowest access latency leaks the confidential byte $k$ with 100\% fidelity.

\section{SiliconTree System Architecture}

\subsection{Reverse-Order Pipeline Evaluation Paradigm}
In a physical synchronous hardware pipeline, all pipeline stage registers update concurrently on the rising clock edge ($\uparrow \text{CLK}$). In software simulation, sequential execution of stages introduces latch race conditions if executed forward ($\text{IF} \to \text{ID} \to \text{EX} \to \text{MEM} \to \text{WB}$): stage $S_i$ would overwrite its output latch, feeding newly computed data into stage $S_{i+1}$ within the exact same cycle.

To solve this without allocating duplicative intermediate double-buffers, SiliconTree executes stages in \textbf{reverse order}:
\begin{equation}
\text{Evaluate}(\text{ClockTick}) \implies \text{WB}() \to \text{MEM}() \to \text{EX}() \to \text{ID}() \to \text{IF}()
\end{equation}
This guarantees that every stage consumes data produced strictly in cycle $t-1$, achieving hardware-identical latch semantics in $O(1)$ stage complexity.

\subsection{High-Fidelity Arithmetic Logic Unit}
SiliconTree implements algorithmic, cycle-accurate multipliers and dividers:
\begin{itemize}
    \item \textbf{Radix-2 Booth Multiplier:} Computes signed two's complement multiplication by scanning multiplier bits $Q_0$ and $Q_{-1}$ \cite{booth1951}. For bit transitions $01_2$, the multiplicand $M$ is added to accumulator $A$; for $10_2$, $M$ is subtracted; for $00_2$ and $11_2$, no operation occurs. An arithmetic right shift ($[A, Q, Q_{-1}] \gg 1$) is executed per step.
    \item \textbf{Radix-4 Bit-Pair Recoding:} Scans overlapping bit triplets $\{b_{2i+1}, b_{2i}, b_{2i-1}\}$ to generate recoded operation set $\{-2, -1, 0, +1, +2\} \times M$ \cite{macsorley1961}. This halves the required multiplication cycles:
    \begin{equation}
    N_{\text{cycles}}(\text{Radix-4}) = \left\lceil \frac{W}{2} \right\rceil
    \end{equation}
    \item \textbf{Integer Division:} Implements both Restoring Division and Non-Restoring Division, eliminating redundant restore cycles.
\end{itemize}

\subsection{Hierarchical Memory Subsystem}
The memory architecture models physical SRAM caches backed by DRAM:
\begin{itemize}
    \item \textbf{Set-Associative Mapping:} Configurable capacity $C$, line size $L$, and associativity $K$. The number of sets $S$ is given by $S = C / (K \times L)$. Addresses are partitioned into tag, set index, and block offset fields.
    \item \textbf{LRU \& Write-Back:} True LRU replacement maintains access timestamp age counters across ways. A write-back, write-allocate policy tracks dirty bits, issuing DRAM write cycles only upon eviction of modified lines.
    \item \textbf{Microarchitectural Cache Tag State:} Cache lines store an explicit $\texttt{speculative\_touch}$ metadata bit, enabling exact tracking of transient access footprints.
\end{itemize}

\subsection{Speculative Execution and LFENCE Serialization}
When a conditional branch enters the pipeline:
\begin{enumerate}
    \item The branch predictor predicts direction and updates $PC$.
    \item Ensuing instructions entering $\text{IF}$ and $\text{ID}$ are tagged with $\texttt{is\_speculative} = \texttt{True}$.
    \item If an out-of-bounds load executes speculatively, it accesses L1 cache, warms the corresponding set, and tags the line with $\texttt{speculative\_touch} = \texttt{True}$.
    \item When the branch reaches $\text{EX}$: if mispredicted, an immediate flush signal is broadcast, overwriting $\text{IF}$ and $\text{ID}$ latches with bubbles and redirecting $PC$. However, the L1 cache state remains altered.
    \item \textbf{Serialization Mitigation ($\texttt{LFENCE}$):} When $\texttt{LFENCE}$ is decoded in $\text{ID}$, it halts instruction fetch and decode until all preceding instructions retire through $\text{WB}$, fully serializing execution and preventing speculative cache pollution.
\end{enumerate}

\subsection{Immutable State Snapshots and Time-Travel}
Every clock cycle generates an immutable dataclass snapshot (\texttt{CycleSnapshot}) containing global cycle, IPC/CPI metrics, architectural/speculative registers, pipeline stage slots, L1/L2 cache tags, bus packets, and memory events. Snapshots are appended to a bounded deque. Calling $\texttt{clock.rewind}(k)$ restores internal registers deterministically via \texttt{EngineRestoreBundle}, providing instantaneous $O(1)$ time travel.

\section{Experimental Evaluation}

\subsection{Automated Test Suite Verification}
SiliconTree was evaluated through an automated test suite containing 14 unit and integration test gates covering ALU accuracy, hazard bubbles, cache conflict misses, Spectre v1 exploitation, and system performance. All 14 tests pass deterministically in $1.66\,\text{s}$ on Python 3.13.

\subsection{Microarchitectural Attack Characterization (SpectreScope)}
We evaluate the Spectre v1 gadget execution in SiliconTree using an out-of-bounds index ($x=16$) against array bound 4:

\begin{table}[t]
\caption{Microarchitectural Metrics Under Spectre v1 Attack and Mitigation}
\label{tab:spectre}
\centering
\begin{tabular}{lcc}
\toprule
\textbf{Parameter / Metric} & \textbf{Without LFENCE} & \textbf{With LFENCE} \\
\midrule
Total Execution Cycles & 18 cycles & 24 cycles \\
Branch Prediction Result & Mispredicted (EX) & Mispredicted (EX) \\
Speculative Instructions Fetched & 3 & 0 (Serialized) \\
Out-of-Bounds Load Executed & Yes (Transient) & No \\
L1 Lines Warmed Speculatively & 1 (Set $k$) & 0 \\
Flush+Reload Target Recovery & \textbf{100\% Leaked} ($t_{\text{hit}}=1$) & \textbf{0\% Leaked} ($t_{\text{miss}}=9$) \\
Pipeline Stalls / Bubbles & 2 squash cycles & 6 serialization cycles \\
\bottomrule
\end{tabular}
\end{table}

As reported in Table~\ref{tab:spectre}, without $\texttt{LFENCE}$, the processor speculatively executes the gadget, warming Set $k$ in L1 cache before the mispredicted $\texttt{BGE}$ squashes instructions in $\text{EX}$. The simulated Flush+Reload routine reads Set $k$ in 1 cycle versus 9 cycles for non-targeted sets, successfully recovering the confidential byte. When $\texttt{LFENCE}$ is enabled, the pipeline halts fetch at $\text{ID}$ until $\texttt{BGE}$ retires, completely preventing speculative cache pollution.

\subsection{Cache Conflict Misses and Associativity Scaling}
To quantify cache performance, we simulate a strided memory access pattern across alternating addresses mapping to the same cache set ($0\text{x}0200$ and $0\text{x}0280$):

\begin{table}[t]
\caption{Cache Replacement and Conflict Miss Benchmark}
\label{tab:cache}
\centering
\begin{tabular}{lcccc}
\toprule
\textbf{Configuration} & \textbf{Accesses} & \textbf{Hits} & \textbf{Conflict Misses} & \textbf{Hit Rate} \\
\midrule
128B Direct-Mapped ($K=1$) & 20 & 2 & 18 & 10.0\% \\
128B 2-Way Set-Assoc ($K=2$) & 20 & 18 & 2 & 90.0\% \\
256B 4-Way Set-Assoc ($K=4$) & 20 & 20 & 0 & 100.0\% \\
\bottomrule
\end{tabular}
\end{table}

In a direct-mapped configuration ($K=1$), competing addresses repeatedly evict one another, generating 18 conflict misses (Table~\ref{tab:cache}). Upgrading to a 2-way set-associative cache retains both addresses simultaneously, raising the hit rate to 90.0\%, while a 4-way cache achieves zero conflict misses.

\subsection{Comparative Evaluation}
Table~\ref{tab:comparison} contrasts SiliconTree with existing architectural simulation frameworks.

\begin{table}[t]
\caption{Feature Comparison of Computer Architecture Simulators}
\label{tab:comparison}
\centering
\begin{tabular}{lcccc}
\toprule
\textbf{Capability / Feature} & \textbf{gem5 \cite{binkert2011}} & \textbf{MARS \cite{vollmar2006}} & \textbf{SimpleScalar \cite{simplescalar}} & \textbf{SiliconTree} \\
\midrule
Simulation Granularity & Cycle & Functional & Cycle & \textbf{Cycle} \\
5-Stage Hazard \& Forwarding & Yes & No & Yes & \textbf{Yes (Visual)} \\
Set-Associative LRU Cache & Yes & No & Yes & \textbf{Yes (Matrix)} \\
Speculative Execution Modeling & Full O3 & No & Partial & \textbf{Yes (Transient)} \\
Spectre / Flush+Reload Sim & Complex & No & No & \textbf{Yes (1-Click)} \\
Single-Cycle Time-Travel & Checkpoints & No & No & \textbf{Yes ($-1$ Cyc)} \\
Zero-Setup Web Workbench & No & Swing & No & \textbf{Yes (60 FPS)} \\
Learning Curve & Very High & Low & Moderate & \textbf{Low} \\
\bottomrule
\end{tabular}
\end{table}

\section{Conclusion}
In this paper, we presented \textbf{SiliconTree}, an interactive, cycle-accurate computer architecture and hardware security simulation suite. By implementing a strict MVC design pattern with reverse-order pipeline evaluation ($\text{WB} \to \text{MEM} \to \text{EX} \to \text{ID} \to \text{IF}$), SiliconTree resolves pipeline latch overwrites deterministically. The engine faithfully captures microarchitectural side-channel vulnerabilities, including Spectre v1 speculative cache line warming, Flush+Reload timing differentiation, and $\texttt{LFENCE}$ mitigation barriers. Its immutable snapshot deque enables real-time single-cycle time-travel ($-1$ Cyc) and post-mortem forensic crash timeline scrubbing. Evaluated across 14 rigorous test gates and five interactive architectural modes, SiliconTree provides researchers, educators, and system architects with a transparent, highly observable platform for exploring processor performance and hardware security.

\begin{thebibliography}{00}
\bibitem{hennessy2017} J. L. Hennessy and D. A. Patterson, \textit{Computer Architecture: A Quantitative Approach}, 6th ed. Cambridge, MA, USA: Morgan Kaufmann, 2017.
\bibitem{hill2008} M. D. Hill and M. R. Marty, ``Amdahl's law in the multicore era,'' \textit{IEEE Computer}, vol. 41, no. 7, pp. 33--38, Jul. 2008.
\bibitem{ge2021} D. Ge, S. Chen, and K. Basu, ``A survey on hardware-assisted security for microarchitectural attack detection,'' \textit{IEEE Design \& Test}, vol. 38, no. 5, pp. 8--18, Oct. 2021.
\bibitem{chandrasekaran2022} S. T. Chandrasekaran, A. P. Kuruvila, K. Basu, and A. Sanyal, ``Real-time hardware-based malware and micro-architectural attack detection utilizing CMOS reservoir computing,'' \textit{IEEE Trans. Circuits Syst. II, Exp. Briefs}, vol. 69, no. 2, pp. 349--353, Feb. 2022.
\bibitem{kocher2019} P. Kocher \textit{et al.}, ``Spectre attacks: Exploiting speculative execution,'' in \textit{Proc. 40th IEEE Symp. Secur. Privacy (SP)}, San Francisco, CA, USA, May 2019, pp. 1--19.
\bibitem{lipp2018} M. Lipp \textit{et al.}, ``Meltdown: Reading kernel memory from user space,'' in \textit{Proc. 27th USENIX Secur. Symp.}, Baltimore, MD, USA, Aug. 2018, pp. 973--990.
\bibitem{yarom2014} Y. Yarom and K. E. Falkner, ``FLUSH+RELOAD: A high resolution, low footprint practical cache side-channel attack,'' in \textit{Proc. 23rd USENIX Secur. Symp.}, San Diego, CA, USA, Aug. 2014, pp. 719--734.
\bibitem{zombieload2019} M. Schwarz, M. Lipp, D. Moghimi, J. Van Bulck, J. Stecklina, T. Prescher, and D. Gruss, ``ZombieLoad: Cross-privilege-boundary data sampling,'' in \textit{Proc. ACM SIGSAC Conf. Comput. Commun. Secur. (CCS)}, London, UK, Nov. 2019, pp. 753--768.
\bibitem{rowhammer2016} D. Gruss, C. Maurice, and S. Mangard, ``Rowhammer.js: A remote software-induced fault attack in JavaScript,'' in \textit{Proc. 13th Conf. Detection Intrusions Malware Vulnerability Assess. (DIMVA)}, San Sebasti\'an, Spain, Jul. 2016, pp. 300--321.
\bibitem{binkert2011} N. Binkert \textit{et al.}, ``The gem5 simulator,'' \textit{ACM SIGARCH Comput. Archit. News}, vol. 39, no. 2, pp. 1--7, May 2011.
\bibitem{vollmar2006} K. Vollmar and P. Sanderson, ``MARS: An education-oriented MIPS assembly language simulator,'' in \textit{Proc. 37th SIGCSE Tech. Symp. Comput. Sci. Educ.}, Houston, TX, USA, Mar. 2006, pp. 239--243.
\bibitem{rars} F. G. S. L. Rocha, ``RARS: RISC-V Assembler and Runtime Simulator,'' GitHub repository, 2021. [Online]. Available: https://github.com/TheThirdOne/rars
\bibitem{simplescalar} D. Burger and T. M. Austin, ``The SimpleScalar tool set, version 2.0,'' \textit{ACM SIGARCH Comput. Archit. News}, vol. 25, no. 3, pp. 13--25, Jun. 1997.
\bibitem{booth1951} A. D. Booth, ``A signed binary multiplication technique,'' \textit{Q. J. Mech. Appl. Math.}, vol. 4, no. 2, pp. 236--240, 1951.
\bibitem{macsorley1961} O. L. MacSorley, ``High-speed arithmetic in binary computers,'' \textit{Proc. IRE}, vol. 49, no. 1, pp. 67--91, Jan. 1961.
\bibitem{ozsoy2016} M. Ozsoy, K. N. Khasawneh, C. Donovick, I. Gorelik, N. Abu-Ghazaleh, and D. Ponomarev, ``Hardware-based malware detection using low-level architectural features,'' \textit{IEEE Trans. Comput.}, vol. 65, no. 11, pp. 3332--3344, Nov. 2016.
\bibitem{basu2020} K. Basu, P. Krishnamurthy, F. Khorrami, and R. Karri, ``A theoretical study of hardware performance counters-based malware detection,'' \textit{IEEE Trans. Inf. Forensics Secur.}, vol. 15, pp. 512--525, 2020.
\bibitem{briongos2018} S. Briongos, G. Irazoqui, P. Malag\'on, and T. Eisenbarth, ``CacheShield: Detecting cache attacks through self-observation,'' in \textit{Proc. 8th ACM Conf. Data Appl. Secur. Privacy (CODASPY)}, Tempe, AZ, USA, Mar. 2018, pp. 224--235.
\end{thebibliography}

\end{document}
